%! Graphing Polynomial Functions and Modeling — Unit 4, Lesson 4.6 — Homework
\documentclass[10pt]{article}
\usepackage{algebra2-article}
\usepackage{algebra2-boxes}
\newcommand{\TallMath}[1]{$\displaystyle #1\rule[-1.4em]{0pt}{3.2em}$}

\newcommand{\lgrid}[4]{%
  \draw[step=1, linegray!40, very thin] (#1,#3) grid (#2,#4);
  \draw[->, semithick, charcoal] (#1,0)--(#2,0) node[below right, font=\scriptsize]{$x$};
  \draw[->, semithick, charcoal] (0,#3)--(0,#4) node[left, font=\scriptsize]{$y$};}

\begin{document}

\pageheader{Unit 4, Lesson 4.6}{Homework}
\namedateperiod

\vspace{0.08in}

\begin{notesbox}{Practice}
Arms from the degree and leading coefficient. Intercepts from the factors. Cross, touch, or flatten
from the multiplicities. One test value per interval for above or below. Every problem below is those
four moves in some order.

\begin{enumerate}[leftmargin=1.9em, itemsep=8pt]
  \item \textbf{Predict from the equation.} Fill in the table without expanding anything.

    \vspace{3pt}
    \renewcommand{\arraystretch}{1.5}
    \begin{tabularx}{\linewidth}{|X|C{1.2cm}|C{1.4cm}|C{2.4cm}|X|C{1.4cm}|}
    \hline
    \rowcolor{forestbg}
    \textbf{Function} & \textbf{Deg.} & \textbf{Lead} & \textbf{Left / right arm} & \textbf{Zeros, and cross / touch / flatten at each} & \textbf{$y$-int} \\ \hline
    (a) $f(x)=(x-1)(x+4)^2$    & & & & & \\ \hline
    (b) $g(x)=-2(x+2)^3(x-1)$  & & & & & \\ \hline
    (c) $h(x)=x^2(x-5)$        & & & & & \\ \hline
    \end{tabularx}

  \item \textbf{Same zeros, different pictures.} One graph below is $y=x^2(x-3)$ and the other is
    $y=x(x-3)^2$. (Each horizontal gridline is $4$ units.)
    \begin{center}
    \begin{tikzpicture}[scale=0.46]
      \lgrid{-2}{4}{-2}{3}
      \begin{scope}
        \clip (-2,-2) rectangle (4,3);
        \draw[very thick, forest, domain=-1.2:3.6, samples=160, smooth]
          plot (\x,{((\x)*(\x)*(\x) - 3*(\x)*(\x))/4});
      \end{scope}
      \foreach \x in {-1,1,2,3}
        \draw[charcoal] (\x,-0.15)--(\x,0.15) node[above=1pt, font=\tiny]{$\x$};
      \node[font=\scriptsize\bfseries, charcoal] at (-1.4,2.5) {I};
    \end{tikzpicture}
    \hspace{0.4cm}
    \begin{tikzpicture}[scale=0.46]
      \lgrid{-2}{4}{-3}{2}
      \begin{scope}
        \clip (-2,-3) rectangle (4,2);
        \draw[very thick, forest, domain=-0.7:3.8, samples=160, smooth]
          plot (\x,{((\x)*(\x)*(\x) - 6*(\x)*(\x) + 9*(\x))/4});
      \end{scope}
      \foreach \x in {-1,1,2,3}
        \draw[charcoal] (\x,-0.15)--(\x,0.15) node[above=1pt, font=\tiny]{$\x$};
      \node[font=\scriptsize\bfseries, charcoal] at (-1.4,1.5) {II};
    \end{tikzpicture}
    \end{center}

    \vspace{2pt}
    Graph I is \blank{2.6cm} \qquad Graph II is \blank{2.6cm}

    \vspace{2pt}
    What single feature did you use? \blank{6.0cm}

  \item \textbf{The full five steps.} $f(x) = (x+2)(x-1)(x-4)$

    \vspace{2pt}
    (a) Degree \blank{0.8cm}, leading coefficient \blank{0.8cm}, end behavior \blank{2.6cm}

    \vspace{2pt}
    (b) Zeros: \blank{3.0cm} \; Behavior at each: \blank{2.6cm} \; $y$-intercept: \blank{1.6cm}

    \vspace{2pt}
    (c) Sign chart:

    \vspace{2pt}
    \renewcommand{\arraystretch}{1.5}
    \begin{tabularx}{\linewidth}{|C{2.4cm}|C{1.4cm}|C{2.4cm}|C{1.6cm}|X|}
    \hline
    \rowcolor{forestbg}
    \textbf{Interval} & \textbf{Test $x$} & \textbf{$f(\text{test }x)$} & \textbf{Sign} & \textbf{Above / below} \\ \hline
    $x<-2$   & $-3$ & & & \\ \hline
    $-2<x<1$ & $0$  & & & \\ \hline
    $1<x<4$  & $2$  & & & \\ \hline
    $x>4$    & $5$  & & & \\ \hline
    \end{tabularx}

    \vspace{2pt}
    (d) On which intervals is $f(x)>0$? \blank{5.0cm}

  \item \textbf{Factor first.} $f(x) = x^3 + 2x^2 - 8x$

    \vspace{2pt}
    (a) Factored completely: \blank{4.6cm} \quad (b) Zeros: \blank{3.0cm}

    \vspace{2pt}
    (c) $y$-intercept: \blank{1.4cm} \quad (d) At most how many turning points? \blank{1.0cm}

    \vspace{2pt}
    (e) The graph is below the $x$-axis on $0<x<2$. Test $x=1$ to confirm: $f(1) = $ \blank{1.6cm}

  \item \textbf{Read the equation off the graph.} (Each horizontal gridline is $4$ units.)
    \begin{center}
    \begin{tikzpicture}[scale=0.46]
      \lgrid{-3}{4}{-4}{3}
      \begin{scope}
        \clip (-3,-4) rectangle (4,3);
        \draw[very thick, forest, domain=-1.5:2.9, samples=170, smooth]
          plot (\x,{(2*(\x)*(\x)*(\x) - 6*(\x)*(\x) + 8)/4});
      \end{scope}
      \foreach \x in {-2,-1,1,2,3}
        \draw[charcoal] (\x,-0.15)--(\x,0.15) node[below=1pt, font=\tiny]{$\x$};
      \fill[charcoal] (-1,0) circle (2.5pt);
      \fill[charcoal] (2,0)  circle (2.5pt);
    \end{tikzpicture}
    \end{center}

    \vspace{2pt}
    (a) $x$-intercepts \blank{2.4cm}; the touch is at $x=$ \blank{1.0cm}; $y$-intercept
    \blank{1.4cm}

    \vspace{2pt}
    (b) Skeleton: $f(x)=a\,($\blank{1.6cm}$)($\blank{1.6cm}$)^2$ \quad Solve: $a=$ \blank{1.0cm}

    \vspace{2pt}
    (c) Equation: \blank{4.6cm} \quad (d) Why can $a$ \emph{not} be found from an $x$-intercept?
    \blank{4.0cm}

  \item \textbf{All the characteristics.} Here is $p(x) = x^4 - 4x^2 = x^2(x-2)(x+2)$.
    (Each horizontal gridline is $2$ units.)
    \begin{center}
    \begin{tikzpicture}[scale=0.46]
      \lgrid{-3}{3}{-3}{4}
      \begin{scope}
        \clip (-3,-3) rectangle (3,4);
        \draw[very thick, forest, domain=-2.3:2.3, samples=170, smooth]
          plot (\x,{((\x)*(\x)*(\x)*(\x) - 4*(\x)*(\x))/2});
      \end{scope}
      \foreach \x in {-2,-1,1,2}
        \draw[charcoal] (\x,-0.15)--(\x,0.15) node[above=1pt, font=\tiny]{$\x$};
      \foreach \q in {-2,0,2} \fill[charcoal] (\q,0) circle (2.5pt);
    \end{tikzpicture}
    \end{center}
    \begin{enumerate}[label=(\alph*), leftmargin=1.8em, itemsep=4pt]
      \item Domain: \blank{2.6cm} \quad Range: \blank{2.6cm}
      \item Zeros: \blank{2.6cm} \; At $x=0$ the graph \blank{1.8cm} the axis, because that zero has
        multiplicity \blank{0.9cm}.
      \item Relative maximum at \blank{1.6cm}; absolute minimum value about \blank{1.2cm}, at
        $x\approx$ \blank{1.8cm}.
      \item Is the relative maximum also an absolute maximum? \blank{1.0cm} \; Why?
        \blank{4.4cm}
      \item $p(x)<0$ on \blank{4.6cm} \quad (careful at $x=0$)
    \end{enumerate}

  \item \textbf{Modeling.} A $12$-inch by $12$-inch square of sheet metal has a square of side $x$
    cut from each corner; the sides fold up into an open box.
    \begin{enumerate}[label=(\alph*), leftmargin=1.8em, itemsep=4pt]
      \item $V(x) = $ \blank{4.4cm} \quad Expanded: $V(x)=4x^3-48x^2+144x$. Degree \blank{0.8cm}
      \item Feasible domain: \blank{2.4cm} \; Why is $x=7$ impossible? \blank{4.4cm}
      \item Complete the table.

        \vspace{2pt}
        \renewcommand{\arraystretch}{1.4}
        \begin{tabular}{|C{1.4cm}|C{1.4cm}|C{1.4cm}|C{1.4cm}|C{1.4cm}|C{1.4cm}|}
        \hline
        \rowcolor{forestbg}
        $x$ & $1$ & $2$ & $3$ & $4$ & $5$ \\ \hline
        $V(x)$ & $100$ & & $108$ & & $20$ \\ \hline
        \end{tabular}
      \item The largest volume is \blank{1.4cm} cubic inches, from a cut of \blank{1.0cm} inches.
        Is that a relative maximum, an absolute maximum on the feasible domain, or both?
        \blank{3.0cm}
      \item Interpret $V(4)$ in a sentence about boxes. \blank{5.6cm}
    \end{enumerate}

  \item \textbf{Explain.} One or two sentences each.
    \begin{enumerate}[label=(\alph*), leftmargin=1.8em, itemsep=4pt]
      \item A classmate graphs $y=(x-5)^2(x+2)$ and reports that the graph is above the $x$-axis for
        $x>-2$ except at $x=5$, where it touches. Then they say ``so the graph changes from positive
        to negative at $x=5$.'' Fix the sentence and explain the rule they missed. \writelines{2}
      \item Why can a graph never tell you the exact degree of a polynomial --- only a minimum?
        Give both reasons you met this unit. \writelines{2}
    \end{enumerate}
\end{enumerate}
\end{notesbox}

\begin{extensionbox}
\begin{enumerate}[label=(\alph*), leftmargin=1.8em, itemsep=5pt]
  \item \textbf{Build to order.} Write a polynomial in factored form whose graph touches the
        $x$-axis at $x=3$, crosses at $x=-2$ and $x=0$, and has both arms pointing \emph{down}.
        State its degree and its $y$-intercept, and explain why no smaller degree could work.
  \item \textbf{Impossible?} Could a polynomial graph have exactly $4$ $x$-intercepts and only
        $2$ turning points? Decide, and justify with the turning-point rule.
  \item \textbf{Minimum degree, then count.} A polynomial with real coefficients has a graph that
        crosses the $x$-axis at $x=-1$ and $x=4$, touches it at $x=2$, and has $3$ turning points.
        What is its least possible degree? If the degree is exactly that, how many imaginary zeros
        does it have (Lesson 4.5)?
  \item \textbf{Two boxes, one volume.} Using $V(x)=4x^3-48x^2+144x$ from problem 7, show that a
        cut of $x=3$ and some other feasible cut give the same volume of $108$ cubic inches. (Solve
        $4x^3-48x^2+144x=108$: divide by $4$, then use the Rational Root Theorem to find $x=3$ and
        factor the rest. One of the remaining roots is feasible; give it to two decimal places.)
\end{enumerate}
\writelines{3}
\end{extensionbox}

\begin{spiralbox}
\textbf{Where the unit ends.} You started Unit 4 by reading polynomial graphs you were handed
(Lesson 4.0). You then learned to multiply and factor them (4.1--4.2), divide them (4.3), find every
rational zero (4.4), and count the zeros a graph can never show you (4.5). Today all of it collapsed
into one procedure that runs both directions: equation $\rightarrow$ graph, and graph $\rightarrow$
equation. \textbf{Next: the Unit 4 test}, and after it, \textbf{Unit 5: Rational Functions}, where a
polynomial appears in a denominator for the first time --- and where dividing by zero puts brand-new
lines on the graph called \emph{asymptotes}.
\end{spiralbox}

\end{document}
